\clearpage
\begin{examplebox}
\textbf{\textcolor{CExample}{例 0：先用几个数字看清规则（入门）}}
\begin{center}\begin{tikzpicture}[font=\small,>=stealth,line join=round]
\draw[gray!40,rounded corners=3pt] (0,0) rectangle (14.2,8.3);\node at (7.1,7.700000000000001) {奇函数：把相加改成同一条曲线上的比较};
\begin{scope}[shift={(4.1,3.9499999999999997)},x=1.2cm,y=0.22999999999999998cm]\draw[->,gray] (-2.5,0)--(2.5,0) node[right]{$x$};\draw[->,gray] (0,-10)--(0,10) node[above]{$y$};\draw[blue,thick,solid,domain=-2.1:2.1,samples=180] plot(\x,{(\x)^3});\fill[orange!85!black] (1,1) circle(2.2pt);\fill[orange!85!black] (2,8) circle(2.2pt);\end{scope}\node[anchor=west,align=left,text width=5.0cm] at (8.4,4.0) {图取 $f(x)=x^3$\\[9pt]$a=2,\ b=-1$\\[9pt]$f(-b)=f(1)=1$\\[9pt]$f(a)+f(b)=8-1=7>0$\\[9pt]比较 $a=2>-b=1$};
\node[text width=12.6cm,align=center] at (7.1,.65) {图中标出的是比较用的 $a$ 和 $-b$；规律依赖奇性与严格单调性。};
\end{tikzpicture}\end{center}



\textbf{奇函数：}先取 $f(x)=x^3$。这是奇函数，且 $x$ 越大，$x^3$ 越大。
\[
\begin{array}{c|c|c}
(a,b)&a+b&f(a)+f(b)\\\hline
(2,-1)&1>0&8-1=7>0\\
(1,-2)&-1<0&1-8=-7<0\\
(2,-2)&0&8-8=0
\end{array}
\]
\textbf{偶函数：}再取 $f(x)=x^2$，看 $a=-3,b=2$。虽然 $-3<2$，但
\[
|-3|=3>2=|2|,\qquad f(-3)=9>4=f(2).
\]
这说明\textbf{奇函数看 $a+b$ 的正负；偶函数看 $|a|$、$|b|$ 谁更大}，而不是先去算函数值。

\medskip

\end{examplebox}
\clearpage
\begin{examplebox}
\textbf{\textcolor{CExample}{例 1：基础题——奇函数递增，解函数值之和}}
\begin{center}\begin{tikzpicture}[font=\small,>=stealth,line join=round]
\draw[gray!40,rounded corners=3pt] (0,0) rectangle (14.2,8.8);\node at (7.1,8.200000000000001) {例1：奇函数递增，三倍输入之和决定答案};
\node[anchor=west] at (.6,6.8) {输入和 $s=3x-1$，比较 $s$ 与 $0$：};\draw[->,gray] (1,6.3)--(13,6.3);\draw[blue,very thick] (6,6.3)--(12,6.3);\draw[blue,fill=white] (6,6.3) circle(2.7pt);\node[below=7pt] at (6,6.3) {$0$};\node[anchor=west] at (.6,5.0) {换回 $x$，分界 $x=\frac13$：};\draw[->,gray] (1,4.5)--(13,4.5);\draw[violet,very thick] (6,4.5)--(12,4.5);\draw[violet,fill=white] (6,4.5) circle(2.7pt);\node[below=7pt] at (6,4.5) {$\frac13$};
\node[text width=12.6cm,align=center] at (7.1,.65) {两行分别表示 $s$ 和 $x$，不是同一数轴尺度；严格大于不含分界点。};
\end{tikzpicture}\end{center}



\textbf{题目：}已知 $f(x)$ 是定义在 $\mathbb R$ 上的奇函数，在 $[0,+\infty)$ 上严格递增，解
\[
 f(x-1)+f(2x)>0.
\]
\textbf{第一步：}奇函数，且右半边严格递增，所以整个实数范围都严格递增。

\textbf{第二步：}把 $a=x-1$、$b=2x$ 代入结论：
\[
 f(x-1)+f(2x)>0\iff (x-1)+2x>0.
\]
\textbf{第三步：}解 $3x-1>0$，得
\[
\boxed{x>\frac13}.
\]
\textbf{提醒：}本题完全不需要知道 $f(x)$ 的具体式子。

\medskip

\end{examplebox}
\clearpage
\begin{examplebox}
\textbf{\textcolor{CExample}{例 2：偶函数递增，求参数范围}}
\begin{center}\begin{tikzpicture}[font=\small,>=stealth,line join=round]
\draw[gray!40,rounded corners=3pt] (0,0) rectangle (14.2,8.8);\node at (7.1,8.200000000000001) {例2：绝对值大于1，两边的解都要保留};
\node[anchor=west] at (.6,6.8) {令 $u=2a-1$，解 $|u|>1$：};\draw[->,gray] (1,6.3)--(13,6.3);\draw[blue,very thick] (1.3,6.3)--(5,6.3);\draw[blue,very thick] (8,6.3)--(12,6.3);\draw[blue,fill=white] (5,6.3) circle(2.7pt);\node[below=7pt] at (5,6.3) {$-1$};\draw[blue,fill=white] (8,6.3) circle(2.7pt);\node[below=7pt] at (8,6.3) {$1$};\node[anchor=west] at (.6,5.0) {换回 $a$，解为 $a<0$ 或 $a>1$：};\draw[->,gray] (1,4.5)--(13,4.5);\draw[violet,very thick] (1.3,4.5)--(5,4.5);\draw[violet,very thick] (8,4.5)--(12,4.5);\draw[violet,fill=white] (5,4.5) circle(2.7pt);\node[below=7pt] at (5,4.5) {$0$};\draw[violet,fill=white] (8,4.5) circle(2.7pt);\node[below=7pt] at (8,4.5) {$1$};
\node[text width=12.6cm,align=center] at (7.1,.65) {各行是各自的数轴；空心端点表示不取等号，两条射线缺一不可。};
\end{tikzpicture}\end{center}



\textbf{题目：}已知 $f(x)$ 是定义在 $\mathbb R$ 上的偶函数，在 $[0,+\infty)$ 上严格递增，且
\[
 f(1)<f(2a-1).
\]
求 $a$ 的范围。

\textbf{第一步：}将两个输入变成到 $0$ 的距离：$|1|$ 与 $|2a-1|$。

\textbf{第二步：}右半边严格递增，因此
\[
 f(1)<f(2a-1)\iff 1<|2a-1|.
\]
\textbf{第三步：}一个数的绝对值大于 $1$，意味着它小于 $-1$ 或大于 $1$：
\[
 2a-1<-1\quad\text{或}\quad 2a-1>1.
\]
于是
\[
\boxed{a<0\quad\text{或}\quad a>1}.
\]
\textbf{数轴理解：}$2a-1$ 必须落在 $[-1,1]$ 的外面，两边都要保留。

\medskip

\end{examplebox}
\clearpage
\begin{examplebox}
\textbf{\textcolor{CExample}{例 3：奇函数递减，而且定义域有限}}
\begin{center}\begin{tikzpicture}[font=\small,>=stealth,line join=round]
\draw[gray!40,rounded corners=3pt] (0,0) rectangle (14.2,8.8);\node at (7.1,8.200000000000001) {例3：先求代数解，再保留可以代入的输入};
\node[anchor=west] at (.6,6.8) {比较得到 $t<-\frac12$：};\draw[->,gray] (1,6.3)--(13,6.3);\draw[blue,very thick] (1.3,6.3)--(6.5,6.3);\draw[blue,fill=white] (6.5,6.3) circle(2.7pt);\node[below=7pt] at (6.5,6.3) {$-\frac12$};\node[anchor=west] at (.6,5.0) {输入限制 $-3\le t\le2$：};\draw[->,gray] (1,4.5)--(13,4.5);\draw[violet,very thick] (3,4.5)--(10,4.5);\draw[violet,fill=violet] (3,4.5) circle(2.7pt);\node[below=7pt] at (3,4.5) {$-3$};\draw[violet,fill=violet] (10,4.5) circle(2.7pt);\node[below=7pt] at (10,4.5) {$2$};\node[anchor=west] at (.6,3.1999999999999997) {共同部分 $[-3,-\frac12)$：};\draw[->,gray] (1,2.6999999999999997)--(13,2.6999999999999997);\draw[orange!85!black,very thick] (3,2.6999999999999997)--(6.5,2.6999999999999997);\draw[orange!85!black,fill=orange!85!black] (3,2.6999999999999997) circle(2.7pt);\node[below=7pt] at (3,2.6999999999999997) {$-3$};\draw[orange!85!black,fill=white] (6.5,2.6999999999999997) circle(2.7pt);\node[below=7pt] at (6.5,2.6999999999999997) {$-\frac12$};
\node[text width=12.6cm,align=center] at (7.1,.65) {三行共用同一 $t$ 尺度；最左端可取，右端使和等于0而排除。};
\end{tikzpicture}\end{center}



\textbf{题目：}已知 $f(x)$ 是定义在 $[-3,3]$ 上的奇函数，在 $[0,3]$ 上严格递减，解
\[
 f(t)+f(t+1)>0.
\]
\textbf{第一步：先看方向。}右半边严格递减的奇函数，在整个 $[-3,3]$ 上也严格递减，所以
\[
 f(t)+f(t+1)>0\iff t+(t+1)<0.
\]
得到 $2t+1<0$，即 $t<-\frac12$。

\textbf{第二步：别忘记输入必须能代入函数。}
\[
 -3\leqslant t\leqslant3,\qquad -3\leqslant t+1\leqslant3.
\]
两者同时成立给出 $-3\leqslant t\leqslant2$。

\textbf{第三步：}与 $t<-\frac12$ 取共同部分，得
\[
\boxed{t\in[-3,-\tfrac12)}.
\]
\textbf{检查：}若只写 $t<-\frac12$，会把例如 $t=-10$ 这种根本不能代入 $f(t)$ 的数也算进去。

\medskip

\end{examplebox}
\clearpage
\begin{examplebox}
\textbf{\textcolor{CExample}{例 4：偶函数递减，不能照搬递增的方向}}
\begin{center}\begin{tikzpicture}[font=\small,>=stealth,line join=round]
\draw[gray!40,rounded corners=3pt] (0,0) rectangle (14.2,8.3);\node at (7.1,7.700000000000001) {例4：偶函数递减，只需比较两个输入距离};
\begin{scope}[shift={(2.7,2.161111111111111)},x=0.8cm,y=0.5111111111111111cm]\draw[->,gray] (-2,0)--(5.5,0) node[right]{$x$};\draw[->,gray] (0,-1)--(0,8) node[above]{$y$};\draw[blue,thick,solid,domain=-1.5:5,samples=180] plot(\x,{abs(\x+1)});\draw[violet,thick,dashed,domain=-1.5:5,samples=180] plot(\x,{abs(2*\x-3)});\fill[orange!85!black] (0.6666666666666666,1.6666666666666667) circle(2.2pt);\fill[orange!85!black] (4,5) circle(2.2pt);\end{scope}\node[anchor=west,align=left,text width=5.0cm] at (8.4,4.0) {蓝：$|x+1|$\\[9pt]紫：$|2x-3|$\\[9pt]递减时要求蓝低于紫\\[9pt]交点 $x=\frac23,\ 4$\\[9pt]$x<\frac23$ 或 $x>4$};
\node[text width=12.6cm,align=center] at (7.1,.65) {纵轴是输入距离，不是原函数值；在两交点外，蓝线低于紫线。};
\end{tikzpicture}\end{center}



\textbf{题目：}已知 $f(x)$ 是定义在 $\mathbb R$ 上的偶函数，在 $[0,+\infty)$ 上严格递减，解
\[
 f(x+1)>f(2x-3).
\]
\textbf{第一步：}离 $0$ \emph{较近}的输入，函数值反而较大：
\[
 f(x+1)>f(2x-3)\iff |x+1|<|2x-3|.
\]
\textbf{第二步：}两边都是非负数，可同时平方：
\[
 (x+1)^2<(2x-3)^2.
\]
移项后是
\[
 0<3x^2-14x+8=(3x-2)(x-4).
\]
\textbf{第三步：}两数相乘为正，要么同正，要么同负，得
\[
\boxed{x<\frac23\quad\text{或}\quad x>4}.
\]
\textbf{代值自检：}取 $f(x)=-x^2$（偶且右半边严格递减），$x=0$ 时题目变成 $-1>-9$，成立；答案也包含 $0$。

\medskip

\end{examplebox}
\clearpage
\begin{examplebox}
\textbf{\textcolor{CExample}{例 5：图像性质给出的“最小值”}}
\begin{center}\begin{tikzpicture}[font=\small,>=stealth,line join=round]
\draw[gray!40,rounded corners=3pt] (0,0) rectangle (14.2,8.3);\node at (7.1,7.700000000000001) {偶函数的右半边递增：最小值在中间};
\begin{scope}[shift={(4.1,2.0951612903225807)},x=0.5263157894736842cm,y=0.14838709677419354cm]\draw[->,gray] (-5.7,0)--(5.7,0) node[right]{$x$};\draw[->,gray] (0,-3)--(0,28) node[above]{$y$};\draw[blue,thick,solid,domain=-5:5,samples=180] plot(\x,{(\x)^2});\fill[orange!85!black] (-4,16) circle(2.2pt);\fill[orange!85!black] (-2,4) circle(2.2pt);\fill[orange!85!black] (0,0) circle(2.2pt);\fill[orange!85!black] (2,4) circle(2.2pt);\end{scope}\node[anchor=west,align=left,text width=5.0cm] at (8.4,4.0) {用 $f(x)=x^2$ 示范\\[9pt]定义域 $[-5,5]$\\[9pt]$f(-4)=16,\ f(2)=4$\\[9pt]$f(0)=0$\\[9pt]$f(-2)=f(2)$\\[9pt]一般结论：$f(0)$ 最小};
\node[text width=12.6cm,align=center] at (7.1,.65) {最小值所在位置由对称和右半边递增决定；具体最小值不一定是0。};
\end{tikzpicture}\end{center}



\textbf{题目：}已知 $f$ 是定义在 $[-5,5]$ 上的偶函数，在 $[0,5]$ 上严格递增。问 $f(-4)$、$f(0)$、$f(2)$ 谁大？什么时候 $f(x)=f(2)$？

\textbf{比较大小：}三者离 $0$ 的距离为 $4,0,2$，所以
\[
\boxed{f(-4)>f(2)>f(0)}.
\]
\textbf{求相等位置：}因为 $f(x)=f(2)\iff |x|=2$，所以
\[
\boxed{x=-2\quad\text{或}\quad x=2}.
\]
\textbf{亮点：}偶函数右半边严格递增时，$f(0)$ 就是整个定义域内的最小值，直接从图像对称性看出来。

\medskip

\end{examplebox}
\clearpage
\begin{examplebox}
\textbf{\textcolor{CExample}{例 6：综合题——绝对值还要和定义域取共同部分}}
\begin{center}\begin{tikzpicture}[font=\small,>=stealth,line join=round]
\draw[gray!40,rounded corners=3pt] (0,0) rectangle (14.2,8.8);\node at (7.1,8.200000000000001) {例6：绝对值比较后，还要与两项定义域取交集};
\node[anchor=west] at (.6,6.8) {代数解 $x<0$ 或 $x>2$：};\draw[->,gray] (1,6.3)--(13,6.3);\draw[blue,very thick] (1.4,6.3)--(5,6.3);\draw[blue,very thick] (9,6.3)--(12,6.3);\draw[blue,fill=white] (5,6.3) circle(2.7pt);\node[below=7pt] at (5,6.3) {$0$};\draw[blue,fill=white] (9,6.3) circle(2.7pt);\node[below=7pt] at (9,6.3) {$2$};\node[anchor=west] at (.6,5.0) {两个输入都可代入：$-\frac12\le x\le1$：};\draw[->,gray] (1,4.5)--(13,4.5);\draw[violet,very thick] (4,4.5)--(7,4.5);\draw[violet,fill=violet] (4,4.5) circle(2.7pt);\node[below=7pt] at (4,4.5) {$-\frac12$};\draw[violet,fill=violet] (7,4.5) circle(2.7pt);\node[below=7pt] at (7,4.5) {$1$};\node[anchor=west] at (.6,3.1999999999999997) {最终 $[-\frac12,0)$：};\draw[->,gray] (1,2.6999999999999997)--(13,2.6999999999999997);\draw[orange!85!black,very thick] (4,2.6999999999999997)--(5,2.6999999999999997);\draw[orange!85!black,fill=orange!85!black] (4,2.6999999999999997) circle(2.7pt);\node[below=7pt] at (4,2.6999999999999997) {$-\frac12$};\draw[orange!85!black,fill=white] (5,2.6999999999999997) circle(2.7pt);\node[below=7pt] at (5,2.6999999999999997) {$0$};
\node[text width=12.6cm,align=center] at (7.1,.65) {三行共用同一 $x$ 尺度；右侧 $x>2$ 整段被定义域限制排除。};
\end{tikzpicture}\end{center}



\textbf{题目：}偶函数 $f$ 定义在 $[-2,2]$，在 $[0,2]$ 上严格递增，解
\[
 f(2x-1)>f(x+1).
\]
\textbf{先去掉 $f$：}
\[
 |2x-1|>|x+1|\iff (2x-1)^2>(x+1)^2.
\]
整理得
\[
 3x(x-2)>0\iff x<0\ \text{或}\ x>2.
\]
\textbf{再查两个输入：}
\[
 -2\leqslant2x-1\leqslant2\iff-\tfrac12\leqslant x\leqslant\tfrac32,
\]
\[
 -2\leqslant x+1\leqslant2\iff-3\leqslant x\leqslant1.
\]
共同部分为 $[-\frac12,1]$，最后
\[
\boxed{x\in[-\tfrac12,0)}.
\]
\textbf{关键：}即使成功把函数不等式变成普通不等式，也不能忘记原来函数能不能取到那些输入。
\end{examplebox}
